\documentclass[../main.tex]{subfiles}

\begin{document}

\chapter{Solutions for Practice Problems}

Here are the solutions to the practice problems at the end of each
notes.

\section{Solutions for Note 1}
\begin{enumerate}
\item
\end{enumerate}

\section{Solutions for Note 2}
\begin{enumerate}
\item
\begin{enumerate}
\item
Consider the following computation.
\[
A - 2B =
\begin{bmatrix}
    1 & 2 \\
    2 & 0
\end{bmatrix}
-
2\begin{bmatrix}
    -1 & 2 \\
    0  & 3
\end{bmatrix}
=
\begin{bmatrix}
    1 & 2 \\
    2 & 0
\end{bmatrix}
+
\begin{bmatrix}
    2 & -4 \\
    0 & -6
\end{bmatrix}
=
\begin{bmatrix}
    3 & -2 \\
    2 & -6
\end{bmatrix}
\]
\item
Consider the following equation.
\[
AB =
\begin{bmatrix}
    1 & 2 \\
    2 & 0
\end{bmatrix}
\begin{bmatrix}
    -1 & 2 \\
    0  & 3
\end{bmatrix}
=
\begin{bmatrix}
    -1 & 8 \\
    -2 & 7
\end{bmatrix}
\]
\item
Because $AB$ is found above, $BA$ could be found.
\[
BA =
\begin{bmatrix}
    -1 & 2 \\
    0  & 3
\end{bmatrix}
\begin{bmatrix}
    1 & 2 \\
    2 & 0
\end{bmatrix}
=
\begin{bmatrix}
    3 & -2 \\
    6 & 0
\end{bmatrix}
\]
Therefore, the following equation is true.
\[
(AB - BA)^\intercal =
\left(
\begin{bmatrix}
    -1 & 8 \\
    -2 & 7
\end{bmatrix}
-
\begin{bmatrix}
    3 & -2 \\
    6 & 0
\end{bmatrix}
\right)^\intercal
=
\left(
\begin{bmatrix}
    -4 & 10 \\
    -8 & 7
\end{bmatrix}
\right)^\intercal
=
\begin{bmatrix}
    -4 & -8 \\
    10 & 7
\end{bmatrix}
\]
\end{enumerate}

\item
By the definition of matrix multiplication, $a_2B = [4, 5, 6]$.

\item
Notice that $A$ is a $2 \times 2$ square matrix. Therefore, define
$A$ as the following.
\[
A =
\begin{bmatrix}
    a_{11} & a_{12} \\
    a_{21} & a_{22}
\end{bmatrix}
\]
Therefore, the following equation is obtained.
\begin{align*}
2A - 3A^\intercal
&=
\begin{bmatrix}
    2a_{11} & 2a_{12} \\
    2a_{21} & 2a_{22}
\end{bmatrix}
-
\begin{bmatrix}
    3a_{11} & 3a_{21} \\
    3a_{12} & 3a_{22}
\end{bmatrix}
=
\begin{bmatrix}
    2a_{11} - 3a_{11} & 2a_{12} - 3a_{21} \\
    2a_{21} - 3a_{12} & 2a_{22} - 3a_{22}
\end{bmatrix} \\
&=
\begin{bmatrix}
    -1 & -5 \\
    0  & -4
\end{bmatrix}
\end{align*}
Continuing, the elements of $A$ could be found.
\begin{align*}
    2a_{11} - 3a_{11} &= -1 \\
    2a_{12} - 3a_{21} &= -5 \\
    2a_{21} - 3a_{12} &= 0  \\
    2a_{22} - 3a_{22} &= -4
\end{align*}
Therefore, $a_{11} = 1$, $a_{12} = 2$, $a_{21} = 3$, and $a_{22} = 4$.
Thus, $A$ can be defined as the following.
\[
A =
\begin{bmatrix}
    1 & 2 \\
    3 & 4
\end{bmatrix}
\]

\item
For any square matrices $S$, define $A$ and $B$ as the following.
\[
A = \frac{S + S^\intercal}{2}, \quad
B = \frac{S - S^\intercal}{2}
\]
Notice that $S = A + B$. Therefore, it suffices to show that $A$ is
symmetric and $B$ is skew-symmetric. By Theorem \ref{Properties of
Transpose of a Matrix}, the following equations hold.
\[
A^\intercal
= \frac{S + S^\intercal}{2}, \quad
B^\intercal
= \frac{S^\intercal - S}{2}
= -B
\]
Thus, all square matrices can be represented as the sum of a symmetric
and skew-symmetric matrices.

\item
By definition, there exists $n, m \in \mathbb{Z}$ such that
$A^n = 0$ and $B^m = 0$. Because $AB = BA$, $A^{nm} B^{nm} = 0$ can
be written as $(AB)^{nm} = 0$ and $AB$ is nilpotent.

\item
By Theorem \ref{Properties of Transpose of a Matrix}, $(AB +
(AB)^\intercal)^\intercal = (AB)^\intercal + AB = AB +
(AB)^\intercal$. Therefore, $AB + (AB)^\intercal$ is symmetric.
\end{enumerate}

\section{Solutions for Note 3}
\begin{enumerate}
\item
\end{enumerate}
\end{document}


